\chapter{Opportunity Identification}
\label{ch:opportunity-identification}

\section{Session Overview}

This session turns the innovation charter from Session~\ref{sec:innovation-charter} into a set of concrete opportunity options. An \textbf{innovation opportunity} is a possible direction for value creation that fits the scope of the charter and is worth exploring further. At this stage, your goal is not to prove that one idea is the final answer. Your goal is to generate several plausible directions and then screen them systematically.

By the end of this session, you should be able to:
\begin{enumerate}
    \item generate several innovation opportunities that align with your charter;
    \item explain how each opportunity could create value for hydroponic users;
    \item evaluate opportunities using consistent screening criteria; and
    \item prepare a prioritised shortlist for later planning work.
\end{enumerate}

\section{Key Terms}

This section introduces terms used throughout the activity.

\begin{itemize}
    \item \textbf{Innovation charter}: the guiding statement created in Session~\ref{sec:innovation-charter} that sets the scope and boundaries of the project.
    \item \textbf{Innovation opportunity}: a possible product or improvement direction that fits the charter.
    \item \textbf{Value creation}: the benefit an opportunity could provide to users, operators, maintainers, or the business context.
\end{itemize}

\section{Pre-Class Activity: Generate Opportunity Options}

Before class, review your team's final innovation charter and generate at least FIVE opportunity options that fit within its scope.

\subsection{How to Brainstorm}

% When generating ideas, consider:
% \begin{itemize}
%     \item user frustrations or unmet needs in hydroponic systems;
%     \item technical issues involving flow, sensing, structure, maintenance, or energy use;
%     \item improvements to safety, efficiency, reliability, or usability; and
%     \item emerging technologies that could support a better solution.
% \end{itemize}

% Try to keep the opportunities broad enough for discussion but specific enough to describe clearly. For example, ``improve nutrient monitoring for small hydroponic growers'' is a stronger opportunity statement than ``make hydroponics better.''
\begin{enumerate}
    \item Refer to the innovation charter you developed in Session~\ref{sec:innovation-charter} as a guide for generating new opportunities.
    \item Generate a minimum of five opportunities for innovation that align with the objectives and boundaries outlined in your innovation charter.
    \begin{enumerate}
        \item Try to think outside the box (without the help of any LLM) and explore multitude possibilities.
        \item Consider various aspects including market needs, technological advancements, customer preferences, or process improvements.
        \item Most important, ensure that the opportunities you come out should within the confine of the broad scope of the innovation charter and have the potential to bring value and impact to your team project.
    \end{enumerate}
    \item Note-down each opportunity.

    While writing down each opportunity, consider the following optional steps, which are highly recommended to facilitate ease of revision.
    \begin{enumerate}
        \item Include a descriptive explanation of the idea and its relationship to the innovation charter.
        \item Consider the potential benefits, challenges, or risks, and identify any necessary resources or support required to pursue each opportunity.
    \end{enumerate}
\end{enumerate}

\subsection{Recording Format}

Document each opportunity using the template in Table~\ref{tab:session4-opportunity-template}. Adding short notes now will make later screening easier.

\begin{table}[htbp]
\centering
\caption{Template for recording innovation opportunities}
\label{tab:session4-opportunity-template}
\renewcommand{\arraystretch}{1.2}
\begin{tabularx}{\textwidth}{|>{\raggedright\arraybackslash}p{0.06\textwidth}|>{\raggedright\arraybackslash}X|>{\raggedright\arraybackslash}p{0.18\textwidth}|>{\raggedright\arraybackslash}p{0.18\textwidth}|>{\raggedright\arraybackslash}p{0.18\textwidth}|}
\hline
\textbf{No.} & \textbf{Opportunity description} & \textbf{Charter alignment} & \textbf{Expected value} & \textbf{Physics or technical focus} \\
\hline
1 &  &  &  &  \\
\hline
2 &  &  &  &  \\
\hline
3 &  &  &  &  \\
\hline
\end{tabularx}
\end{table}

\section{Example Opportunity}

A charter focused on compact hydroponic support products could lead to an opportunity such as: ``Develop a low-cost sensing accessory that helps small growers detect nutrient-level changes quickly.'' This opportunity aligns with the charter, creates value through easier monitoring, and opens clear applied-physics questions related to instrumentation and calibration.

\section{In-Class Activity: Screen and Prioritise Opportunities}

In class, combine the opportunities generated by each team member and screen them as a group. The objective is not to identify one final winner immediately. Instead, remove weak ideas and keep the strongest opportunities for deeper work.

\subsection{Screening Criteria}

Rate each opportunity using the following criteria:
\begin{enumerate}
    \item \textbf{Alignment}: Does the opportunity fit the innovation charter closely?
    \item \textbf{Impact}: Would it solve a meaningful user or system problem?
    \item \textbf{Feasibility}: Can the team study and develop it within the course constraints?
    \item \textbf{Technical depth}: Does it allow relevant applied physics analysis or justification?
    \item \textbf{Resource fit}: Are the likely materials, data, and skills reasonably accessible?
\end{enumerate}

\subsection{Scoring Template}

Use a 1--5 scale for each criterion, where 5 is the strongest score. Table~\ref{tab:session4-screening-matrix} can be used to structure the discussion.

\begin{table}[htbp]
\centering
\caption{Opportunity screening matrix}
\label{tab:session4-screening-matrix}
\renewcommand{\arraystretch}{1.2}
\begin{tabularx}{\textwidth}{|>{\raggedright\arraybackslash}X|c|c|c|c|c|>{\raggedright\arraybackslash}p{0.16\textwidth}|}
\hline
\textbf{Opportunity} & \textbf{Align.} & \textbf{Impact} & \textbf{Feas.} & \textbf{Tech.} & \textbf{Res.} & \textbf{Notes or rationale} \\
\hline
 &  &  &  &  &  &  \\
\hline
 &  &  &  &  &  &  \\
\hline
 &  &  &  &  &  &  \\
\hline
\end{tabularx}
\end{table}

If the lecturer provides voting stickers or another group-decision method, use that process after the scoring discussion. Record the rationale for the final shortlist rather than only the final ranking.

\section{Deliverable and Next Steps}

At the end of the session, your group should have:
\begin{itemize}
    \item a compiled opportunity list;
    \item a completed screening or voting record; and
    \item a prioritised shortlist with brief justifications.
\end{itemize}

Keep these notes because the shortlisted opportunities will feed directly into the next stages of product planning and concept development.
